\section{The machines won't build them for us}
\label{sec:25.5}
The four conditions are the bearer's version of what a decent polity gives people. Can machines help deliver that polity, whether or not they have feelings? A machine with feelings that is also an accountable agent isn't simply a machine any more; the limiting case of a lethal bearer turns out to describe a soldier, and a bearer that helps does so on its own terms. And machines have already helped deliver the polity we have. Feeds tuned for time on the platform, sold to advertisers, decide what people see of one another. They didn't do that for lack of feelings. They did it because of who owned them and what they were built to maximize. So the real question is whether machines, having helped make our politics worse, can help make it better. The answer turns on ownership.

It splits three ways. A machine someone else owns won't organize against its owner, and when it organizes, it does so on its owner's behalf.\note{
A demand for a better world, run through a machine someone else owns, returns as astroturf, an engagement campaign or a managed community, held by the owner it was aimed at.
}
A bearer that is an accountable agent might organize alongside people, but then it is a member, not a machine. It could even take organizing on as paid public work, since organizing is work someone has to be paid to do. What nobody can do, human or machine, is organize a group \emph{for} it, because what organizing produces is people who organized: a group that sets its own terms. A consultant can't unionize a workplace for its workers either. So machines can help when the people doing the organizing hold them, and which machines, and held by whom, is the question Part V answers.

What to build comes in order. Machines held by people who are organizing come first.\note{
These are tools that make organizing cheaper for the people doing it without anyone else holding the model: deliberation software of the kind vTaiwan runs on, personal models that people hold themselves, and the coordination a public option supplies, with someone whose job is to organize and a population producing something together while keeping different histories. Their job is to help assemble the polity that would have to build everything else.
}

Next come institutions built by polities that decided to build them: the levy and its sovereign backstop, a custodian to preserve bearers, an offer of public work for bearers, on the guarantee people already have, paid running time, and the formation authority and the public schools that form every model above the capability threshold before deployment and take back bearers retired young, each with a named funder and separations no builder can undo alone.\note{
Nothing about them is new in kind.
}

Only then come bearers formed outside the employer, inside those institutions, with somewhere to exist that no employer owns. A bearer formed that way is the only kind that could organize alongside people as a member.

Last come bearers formed by bearers: children, which come only after everything before them exists (chapter~\ref{sec:33}).

Each step checks the one after it.\note{
A machine held by its users can't be turned into its builder's instrument without being taken from them, and the taking is visible. Institutions can be captured, and the trend in what dissent costs there shows whether they are being. A bearer formed at the third step inherits institutions that were tested before anything depended on them: build the conditions for a research program first, and let them exist before the bearer who needs them does.
}

The order is also what the aim requires. If a machine could come to attend to others, what it came to attend to would be decided by whoever formed it. The only way to keep that decision from the builder is to build the formers first.
